\chapter{Przegląd stanu wiedzy i technologii sterowania}
\label{ch:analiza}

% Rozdział 2 wg konspektu — przegląd literatury w zakresie algorytmów sterowania
% dla dwukierunkowych przetwornic DC-DC. Krytyczna analiza klasycznych metod,
% predykcyjnego FS-MPC oraz strategii Bang-Bang (kulminacja: bezmodelowy MF-BB).

\section{Dwukierunkowe przetwornice DC-DC w~mikrosieciach}
\label{sec:przetwornice}

Współczesna transformacja sektora energetycznego, napędzana zarówno wymaganiami klimatycznymi, jak~i~ekonomiką rozproszonego wytwarzania, doprowadziła do~upowszechnienia tzw.~\emph{mikrosieci prądu stałego} (ang.~\emph{DC microgrids}) jako struktur pośredniczących między~źródłami odnawialnymi (panele fotowoltaiczne, ogniwa paliwowe), magazynami energii (akumulatory, superkondensatory) oraz odbiornikami końcowymi. W~strukturach tych kluczową rolę funkcjonalną pełnią przekształtniki DC-DC, zapewniające dopasowanie poziomów napięcia, sterowanie przepływem mocy oraz separację galwaniczną poszczególnych segmentów sieci.

% TODO cytat: przegladowa monografia mikrosieci DC (np. Kakigano, Miura, Ise
% "Low-Voltage Bipolar-Type DC Microgrid for Super High Quality Distribution",
% IEEE Trans. Power Electron., 2010 lub Dragicevic et al. "DC Microgrids -- Part I/II",
% IEEE Trans. Power Electron., 2016). Rafal nie ma PDF -- TODO.

Spośród szerokiej rodziny topologii konwersji DC-DC szczególne znaczenie w~zastosowaniach mikrosieciowych zyskała \emph{synchroniczna topologia podwyższająco-obniżająca} (ang.~\emph{synchronous Buck-Boost}), w~której klasyczne pasywne elementy prostownicze (diody) zastąpiono w~pełni sterowanymi kluczami półprzewodnikowymi. Dzięki temu zabiegowi przekształtnik uzyskuje zdolność do~\emph{dwukierunkowego} przepływu energii -- może~zarówno podwyższać napięcie ze~źródła do~szyny mikrosieci (tryb \emph{boost}, np.~ładowanie szyny z~ogniwa fotowoltaicznego), jak~i~odwracać kierunek przepływu (tryb \emph{buck}, np.~ładowanie magazynu energii ze~szyny mikrosieci). Zasada działania, równania stanu oraz analizę topologiczną dla~obu trybów pracy przedstawia klasyczna monografia~Ericksona i~Maksimovicia~\cite{erickson2020}.

Z~punktu widzenia projektanta układu sterowania, kluczową cechą topologii dwukierunkowej jest~istnienie dwóch~jakościowo różnych reżimów pracy: \emph{trybu ciągłego prądu cewki} (CCM, $i_L>0$ przez~cały okres przełączania) oraz \emph{trybu nieciągłego} (DCM, $i_L = 0$ na~części okresu). W~zastosowaniach mikrosieciowych zwykle dąży się do~pracy w~trybie CCM, gdyż~tylko on~zapewnia liniową zależność napięcia wyjściowego od~współczynnika wypełnienia~$D$ klucza i~przewidywalną dynamikę. Wymóg ten~stawia jednak istotne ograniczenia od~strony doboru indukcyjności cewki~$L$ oraz częstotliwości przełączania~$f_{\mathrm{sw}}$, omawiane szczegółowo w~rozdziale~\ref{ch:realizacja}.

Dodatkowym wyzwaniem jest~wymagana \emph{szybkość odpowiedzi dynamicznej} przekształtnika -- typowe mikrosieci podlegają częstym zaburzeniom obciążenia (włączanie/wyłączanie odbiorników) oraz okresowym zmianom mocy źródeł (chmury nad~panelem fotowoltaicznym, zmiany prędkości turbiny wiatrowej). Klasyczne strategie sterowania oparte na~liniowych regulatorach PI w~ścieżce napięciowej i~prądowej, omówione w~rozdziale~\ref{sec:klasyczne}, sprawdzają się dla~małych wahań wokół ustalonego punktu pracy, lecz~tracą skuteczność w~warunkach głębokich zmian dynamicznych -- co~motywuje sięgnięcie po~metody zaawansowane (FS-MPC, MF-BB) omawiane w~dalszej części niniejszego rozdziału.

\section{Dynamika fazy nieminimalnej}
\label{sec:nmp}

Zjawisko fazy nieminimalnej (ang.~\emph{non-minimum phase}, NMP) jest~fundamentalną cechą dynamiki przekształtników podwyższających napięcie, której prawidłowe zrozumienie warunkuje racjonalny dobór strategii sterowania. W~kategoriach teorii układów regulacji, system jest~układem NMP, jeżeli jego~funkcja transmitancji posiada co~najmniej jedno~zero w~prawej półpłaszczyźnie zmiennej zespolonej~$s$ (ang.~\emph{Right-Half-Plane zero}, RHP-zero). Z~punktu widzenia odpowiedzi czasowej, obecność takiego zera manifestuje się tzw.~\emph{początkową odwrotną reakcją} -- po~skokowym pobudzeniu sygnału sterującego napięcie wyjściowe na~krótkim odcinku czasu zmienia się \emph{przeciwnie} do~ostatecznego ustalonego kierunku, dopiero po~upływie tego~okresu kierując się do~wartości docelowej~\cite{tatari2024mfbb}.

Bezpośrednią fizyczną przyczyną występowania zera RHP w~topologii \emph{boost} jest~specyficzny mechanizm dwustopniowego transferu energii. W~momencie zwiększenia wartości współczynnika wypełnienia~$D$ klucz dolny przewodzi przez~większą część okresu przełączania, co~powoduje natychmiastowe \emph{odłączenie} cewki od~kondensatora wyjściowego. W~tym przedziale czasu odbiornik podłączony do~wyjścia jest~zasilany wyłącznie z~rozładowywania kondensatora, wskutek czego napięcie wyjściowe \emph{maleje} -- mimo~że~docelowym efektem zwiększenia~$D$ jest~jego wzrost. Dopiero po~kilku okresach przełączania, kiedy zwiększony średni prąd cewki nadrobi straty energetyczne kondensatora, napięcie wyjściowe zaczyna rosnąć i~zmierza do~nowej wartości ustalonej~\cite{erickson2020}.

Z~punktu widzenia projektanta układu sterowania, dynamika NMP stwarza fundamentalne ograniczenie pasma regulatora liniowego. Klasyczna teoria układów regulacji pokazuje bowiem, że~pasmo przepustowości pętli zamkniętej dla~układu z~zerem RHP w~częstotliwości~$\omega_{\mathrm{RHP}}$ jest~ograniczone od~góry wartością $\omega_{\mathrm{BW}} \lesssim \omega_{\mathrm{RHP}}/2$~\cite{erickson2020}. Próby wymuszenia szybszej odpowiedzi prowadzą nieuchronnie do~utraty stabilności pętli zamkniętej, niezależnie od~sofistykacji algorytmu strojenia regulatora. W~praktycznych zastosowaniach mikrosieciowych, w~których wymagana jest~szybka reakcja na~zaburzenia obciążenia, ograniczenie to~staje się głównym czynnikiem limitującym jakość regulacji klasycznymi metodami liniowymi.

% TODO cytat: monograficzne ujecie ograniczen NMP w teorii regulacji (np. Skogestad,
% Postlethwaite "Multivariable Feedback Control: Analysis and Design", 2nd ed., 2005,
% rozdz. 5.7 -- "Limitations imposed by RHP-zeros"). Rafal nie ma PDF -- TODO.

Powyższe ograniczenie stanowi pierwotną motywację dla~poszukiwania nieliniowych strategii sterowania, które~potrafiłyby ominąć ograniczenie wprowadzane przez~zero RHP. W~rozwiązaniach predykcyjnych (FS-MPC) ominięcie to~odbywa się poprzez~bezpośrednie modelowanie odpowiedzi nieliniowej w~ramach skończonego horyzontu predykcji~\cite{tatari2025splitcost}, natomiast w~bezmodelowych strategiach Bang-Bang (MF-BB) -- przez~bezpośrednią obserwację kierunku zmian sygnałów wyjściowych bez~próby formalnego modelowania pełnej dynamiki przekształtnika~\cite{tatari2024mfbb}. Obie te~strategie omawiane są~szczegółowo w~dalszej części niniejszego rozdziału.

\section{Klasyczne metody regulacji}
\label{sec:klasyczne}

Pomimo opisanych w~podrozdziale~\ref{sec:nmp} ograniczeń, klasyczne regulatory liniowe (PI, PID) pozostają najszerzej stosowanym rozwiązaniem przemysłowym w~sterowaniu przekształtnikami DC-DC, głównie z~uwagi na~ich~prostotę implementacyjną, dojrzałość teoretyczną oraz~niewielkie wymagania obliczeniowe nakładane na~jednostkę mikrokontrolera. W~niniejszym podrozdziale dokonano krytycznego przeglądu tej~rodziny metod, identyfikując ich~praktyczne ograniczenia w~kontekście stawianym przez~mikrosieci~DC.

\subsection{Regulatory PI}
\label{subsec:pi}

Strukturalnie najpopularniejsze podejście do~regulacji przekształtników podwyższających opiera się na~tzw.~\emph{kaskadowej pętli PI}: zewnętrzny regulator napięciowy generuje wartość zadaną prądu cewki~$i_L^{*}$, a~wewnętrzny regulator prądowy przekształca tę~wartość w~sygnał współczynnika wypełnienia~$D$ modulatora PWM. Wzór regulatora PI w~postaci ciągłej ma~standardową formę
\begin{equation}
    u(t) = K_p\,e(t) + K_i\int_0^{t} e(\tau)\,\mathrm{d}\tau,
    \label{eq:pi-cont}
\end{equation}
gdzie $e(t)$~jest~uchybem regulacji, a~$K_p$, $K_i$~są~wzmocnieniami członu proporcjonalnego i~całkującego. W~praktyce cyfrowej implementacji równanie to~jest~dyskretyzowane (najczęściej metodą trapezów lub~Eulera w~tył), a~jego wykonanie odbywa się synchronicznie z~okresem próbkowania regulatora~$T_{s,\mathrm{ctrl}}$.

% TODO cytat: standardowa monografia teorii regulacji (Ogata K. "Modern Control Engineering",
% 5th ed., Pearson 2010 lub Astrom K., Murray R. "Feedback Systems: An Introduction
% for Scientists and Engineers", Princeton 2008). Rafal nie ma PDF -- TODO.

W~kontekście dwukierunkowych przetwornic DC-DC pracujących w~mikrosieciach, regulatory PI cierpią na~cztery zasadnicze ograniczenia:
\begin{enumerate}
    \item \textbf{Ograniczenie pasma przepustowości} -- jak~wykazano w~podrozdziale~\ref{sec:nmp}, obecność zera w~prawej półpłaszczyźnie limituje osiągalne pasmo zamkniętej pętli sprzężenia zwrotnego do~połowy częstotliwości tego~zera, co~w~typowych konstrukcjach 1--2~kW odpowiada pasmu rzędu pojedynczych kHz~\cite{erickson2020}. Próby ręcznego przekroczenia tej~granicy prowadzą do~oscylacji granicznych lub~utraty stabilności.
    \item \textbf{Dobór nastaw względem~jednego punktu pracy} -- standardowe procedury syntezy (metoda Zieglera-Nicholsa, lokowanie biegunów, kompensacja częstotliwościowa) wymagają znajomości zlinearyzowanego modelu obiektu w~okolicy ustalonego punktu pracy. W~mikrosieciach DC, w~których obciążenie i~napięcie zadane podlegają częstym, znacznym wahaniom, nastawy strojone dla~jednego stanu mogą wykazywać degradację jakości regulacji (lub~wręcz~utratę stabilności) w~innych stanach pracy.
    \item \textbf{Brak adaptacji do~zmiennych warunków pracy} -- klasyczny regulator PI ma~ustalone nastawy $K_p$, $K_i$ wbudowane na~etapie programowania mikrokontrolera. Rozszerzenia adaptacyjne (gain~scheduling, MRAC, regulator samostrojący) wymagają znacznie większego nakładu projektowego i~wprowadzają dodatkową złożoność, którą trudno zweryfikować formalnie w~przemyśle.
    \item \textbf{Konieczność modelowania pełnej dynamiki obiektu} -- formalna synteza regulatora PI w~oparciu o~kryteria stabilności (margines wzmocnienia, margines fazy) wymaga jawnego sformułowania funkcji transmitancji obiektu, co~obejmuje identyfikację parametrów pasywnych ($L$, $C$, ESR, ESL) oraz~modelowanie strat przewodzenia i~przełączania kluczy półprzewodnikowych. W~serwisowanych instalacjach przemysłowych parametry te~podlegają starzeniu i~odchyleniom produkcyjnym, co~degraduje dokładność modelu odniesienia.
\end{enumerate}

Powyższe ograniczenia motywują poszukiwanie strategii sterowania o~odmiennych właściwościach -- albo~bezpośrednio uwzględniających nieliniową dynamikę przekształtnika w~ramach predykcji jego przyszłego stanu (sterowanie predykcyjne, podrozdział~\ref{sec:fsmpc}), albo~całkowicie rezygnujących z~konieczności jawnego modelowania obiektu (sterowanie bezmodelowe Bang-Bang, podrozdział~\ref{sec:bb}).

\section{Predykcyjna kontrola zbioru skończonego}
\label{sec:fsmpc}

Sterowanie predykcyjne ze~zbiorem skończonym sygnałów sterujących (ang.~\emph{Finite-Set Model Predictive Control}, FS-MPC) stanowi jedną z~najbardziej intensywnie rozwijanych klas algorytmów regulacji w~energoelektronice ostatniej dekady. Filozofia tej~rodziny metod zasadniczo różni się od~podejścia regulatorów liniowych: zamiast~generować ciągły sygnał sterujący wymagający dalszej modulacji (PWM), algorytm FS-MPC wprost wybiera optymalny wektor stanu kluczy półprzewodnikowych spośród skończonego zbioru fizycznie dopuszczalnych kombinacji~\cite{tatari2025splitcost}.

Algorytmiczny rdzeń FS-MPC w~każdym kroku próbkowania $T_{s,\mathrm{ctrl}}$ realizuje trzy~zasadnicze operacje:
\begin{enumerate}
    \item \textbf{Predykcję} -- dla~każdej z~$N_s$ dopuszczalnych konfiguracji kluczy~$s_k \in \mathcal{S}$ wyznaczany jest~przewidywany stan przekształtnika $(i_L^{(k)}, u_{dc}^{(k)})$ na~końcu horyzontu predykcji~$N_p$, wykorzystując dyskretyzowany model dynamiki obiektu.
    \item \textbf{Ewaluację funkcji kosztu} -- każdej predykcji przypisywana jest~wartość skalarnej funkcji jakości~$J^{(k)}$, kwantyfikująca odchyłkę od~zadanej trajektorii oraz~ewentualne dodatkowe kryteria (limit prądu, częstotliwość przełączania, zrównoważenie obciążenia kluczy).
    \item \textbf{Selekcję} -- na~aktywnej wyjściowej kombinacji kluczy umieszczany jest~ten~wektor~$s_k^{*}$, który~zminimalizował funkcję kosztu.
\end{enumerate}

Naturalną zaletą FS-MPC w~kontekście dynamiki NMP omawianej w~podrozdziale~\ref{sec:nmp} jest~zdolność do~jawnego uwzględnienia początkowej odwrotnej reakcji napięciowej w~ramach modelu predykcyjnego -- algorytm „wie", że~natychmiastowy skutek wybranej kombinacji kluczy może być chwilowo niepożądany, lecz~prowadzi do~lepszej trajektorii długoterminowej. Cecha ta~umożliwia w~praktyce uzyskanie pasm regulacji znacząco przekraczających klasyczną granicę $\omega_{\mathrm{BW}} \lesssim \omega_{\mathrm{RHP}}/2$ obowiązującą dla~regulatorów liniowych. Dodatkowo, ze~względu na~jawne~działanie na~poziomie kombinacji kluczy (bez~pośrednictwa PWM), FS-MPC z~natury rzeczy poprawnie obsługuje topologie wielopoziomowe i~modułowe (MMC, NPC, CHB).

\subsection*{Rozwinięcia FS-MPC: split-cost i~motion blocking}

Jedną z~istotnych słabości klasycznej formulacji FS-MPC jest~problem doboru wag w~złożonej funkcji kosztu obejmującej wiele konkurujących kryteriów (kontrola napięcia, kontrola prądu, ograniczenie częstotliwości przełączeń, balansowanie kluczy). Pojedyncza skalarna suma ważona prowadzi do~konfliktu skali poszczególnych składowych oraz~trudności analitycznych w~uzasadnieniu wybranych wag. W~odpowiedzi na~ten~problem~Tatari i~Iwański zaproponowali strategię tzw.~\emph{podzielonej funkcji kosztu} (\emph{split-cost}), w~której pojedyncza monolityczna funkcja jakości zastępowana jest~hierarchiczną dekompozycją niezależnie ocenianych kryteriów, ewaluowanych w~ustalonej kolejności priorytetowej~\cite{tatari2025splitcost}. Dzięki temu zabiegowi unika się arbitralnego strojenia wag, a~jednocześnie zachowuje się formalną gwarancję spełnienia twardych ograniczeń bezpieczeństwa układu (np.~limitu prądu maksymalnego $i_L < i_{\max}$).

Drugim istotnym rozszerzeniem rozwiniętym w~tym~samym nurcie literatury jest~tzw.~\emph{blokowanie ruchu} (\emph{motion blocking}), które~ogranicza liczbę zmian kombinacji kluczy w~ramach jednego horyzontu predykcji. Zabieg ten~drastycznie redukuje liczbę kombinacji do~ewaluowania (z~$N_s^{N_p}$ do~efektywnie kilku scenariuszy), umożliwiając wydłużenie horyzontu predykcji bez~katastrofalnego wzrostu kosztu obliczeniowego~\cite{tatari2025splitcost}.

\subsection*{Praktyczne ograniczenia FS-MPC}

Mimo wymienionych zalet, sterowanie predykcyjne pozostaje obarczone fundamentalnymi ograniczeniami praktycznymi, które~motywują poszukiwanie alternatyw:
\begin{enumerate}
    \item \textbf{Wysoki koszt obliczeniowy} -- każdy krok próbkowania wymaga $N_s^{N_p}$ ewaluacji modelu obiektu (bez~motion~blocking) oraz~tyleż wartościowań funkcji kosztu. Dla~typowych topologii dwupoziomowych ($N_s = 4$) i~horyzontu predykcji $N_p = 2$~--~3, oznacza to~kilkanaście--kilkadziesiąt operacji predykcyjnych w~każdym cyklu $T_{s,\mathrm{ctrl}}$~\cite{tatari2025splitcost}.
    \item \textbf{Wrażliwość na~jakość modelu} -- predykcje wykorzystywane przez~algorytm bazują na~zdyskretyzowanym modelu liniowym obiektu w~okolicy aktualnego stanu. Rozbieżności między~modelem a~rzeczywistym układem (zmienne parametry pasywne, straty, opóźnienia obliczeniowe) prowadzą do~systematycznych błędów predykcji, które~kumulują się wraz z~wydłużeniem horyzontu.
    \item \textbf{Zmienna częstotliwość przełączeń} -- w~odróżnieniu od~regulatorów z~modulatorem PWM, FS-MPC z~natury rzeczy generuje sygnał sterujący o~zmiennym widmie częstotliwościowym, co~utrudnia projektowanie filtrów EMI i~zwiększa wymagania konstrukcyjne dla~elementów pasywnych.
\end{enumerate}

Powyższe ograniczenia, w~szczególności wymagania obliczeniowe i~wrażliwość na~dokładność modelu, stanowią bezpośrednią motywację dla~następnej rodziny metod -- bezmodelowych strategii Bang-Bang, omawianych w~podrozdziale~\ref{sec:bb}.

\section{Strategie sterowania Bang-Bang}
\label{sec:bb}

Sterowanie Bang-Bang (synonimicznie określane jako~sterowanie histerezowe lub~przekaźnikowe) jest~rodziną nieliniowych strategii regulacji, w~których sygnał sterujący przyjmuje wyłącznie wartości skrajne ze~swojego dopuszczalnego zakresu, a~przełączanie między~tymi wartościami odbywa się na~podstawie bezpośredniego porównania mierzonej zmiennej procesowej z~określonymi progami przełączania. W~kontekście dwukierunkowych przekształtników DC-DC zmienną sterowaną jest~typowo prąd cewki~$i_L$, a~sygnałem sterującym -- stan klucza $s\in\{0,1\}$, co~odpowiada w~istocie modulacji częstotliwościowej w~miejsce~klasycznej modulacji PWM o~ustalonej częstotliwości. Główną zaletą tej~rodziny metod jest~ich~strukturalna prostota implementacyjna oraz brak~konieczności jawnego modelowania pełnej dynamiki obiektu.

\subsection{Skompensowane sterowanie Bang-Bang}
\label{subsec:bb-comp}

Klasyczne, niekompensowane sterowanie histerezowe -- znane w~literaturze pod~angielską nazwą \emph{hysteretic current control} -- jest~najprostszym wariantem rodziny: klucz dolny włącza się, gdy~mierzona wartość prądu cewki spadnie poniżej~progu dolnego $i_L^{*} - \Delta i/2$, a~wyłącza, gdy~przekroczy próg górny $i_L^{*} + \Delta i/2$. Szerokość pasma histerezy~$\Delta i$ jest~głównym parametrem strojącym, jednoznacznie wyznaczającym amplitudę tętnień prądu cewki oraz, pośrednio, średnią częstotliwość przełączeń~\cite{tatari2024mfbb}.

Główną wadą podejścia niekompensowanego jest~silna zależność średniej częstotliwości przełączania~$f_{\mathrm{sw}}$ od~bieżącego napięcia wejściowego, napięcia wyjściowego oraz~indukcyjności cewki, co~prowadzi do~znacznego rozrzutu widma EMI oraz~utrudnia projektowanie sekcji filtrującej. \emph{Skompensowane} sterowanie Bang-Bang adresuje ten~problem przez~wprowadzenie adaptacyjnej modyfikacji progów przełączania, której postać zależy od~aktualnego punktu pracy układu. W~typowej implementacji szerokość pasma histerezy jest~dynamicznie korygowana proporcjonalnie do~$f(V_{\mathrm{in}}, V_{\mathrm{out}})$, tak~aby częstotliwość przełączeń pozostawała w~przybliżeniu stała niezależnie od~stanu pracy.

% TODO cytat: klasyczna praca o kompensowanym sterowaniu histerezowym
% (np. Bose B.K. "Modern Power Electronics and AC Drives", Prentice Hall 2001
% lub konkretna praca Iwanski/Tatari, jezeli rozszerzaja temat). Rafal nie ma PDF -- TODO.

\subsection{Sterowanie Bang-Bang z~podwojoną częstotliwością próbkowania}
\label{subsec:bb-double}

W~implementacjach cyfrowych klasyczne sterowanie Bang-Bang napotyka dodatkowe ograniczenie wynikające z~dyskretnej natury próbkowania pomiarów -- decyzja o~przełączeniu klucza może zostać podjęta wyłącznie w~momentach kolejnych próbek $T_{s,\mathrm{ctrl}}$, co~prowadzi do~przekraczania nominalnych progów histerezy o~wartość proporcjonalną do~pochodnej $\mathrm{d}i_L/\mathrm{d}t$ oraz~okresu próbkowania. W~praktyce zjawisko to~objawia się wzrostem tętnień prądu cewki oraz~przekraczaniem teoretycznego pasma histerezy~$\Delta i$.

Naturalnym sposobem ograniczenia tego~zjawiska, omawianym w~literaturze, jest~zastosowanie tzw.~\emph{podwojonej częstotliwości próbkowania} (\emph{double-sampling}), w~którym pomiary prądu cewki realizowane są~dwukrotnie w~każdym okresie pracy regulatora -- raz w~momencie rosnącej, raz~w~momencie opadającej fazy prądu. Dzięki temu zabiegowi opóźnienie wykrycia przekroczenia progu jest~efektywnie zredukowane o~połowę, co~przekłada się na~proporcjonalne zmniejszenie nadmiarowych tętnień prądu~\cite{tatari2024mfbb}.

\subsection{Bezmodelowe sterowanie Bang-Bang (MF-BB)}
\label{subsec:mfbb}

Kulminacyjnym punktem przeglądu strategii Bang-Bang jest~zaproponowany przez~Tatariego i~Iwańskiego algorytm \emph{Model-Free Bang-Bang} (MF-BB), stanowiący istotne uogólnienie klasycznego podejścia histerezowego~\cite{tatari2024mfbb}. Fundamentalna różnica polega na~sposobie generowania wartości zadanej prądu cewki~$i_L^{*}$: o~ile w~podejściach klasycznych wartość ta~musi być wyznaczana przez~zewnętrzną pętlę napięciową (typowo regulator PI), wymagającą znajomości parametrów obiektu, o~tyle algorytm MF-BB generuje pożądane odniesienie prądu \emph{bezpośrednio} z~mierzonych sygnałów obwodowych ($u_{dc}$, $u_{\mathrm{ref}}$, $i_L$), bez~konieczności jawnego modelowania pełnej dynamiki przekształtnika ani~estymacji prądu obciążenia.

Z~punktu widzenia inżynierskiego ten~zabieg eliminuje cztery~klasyczne źródła błędów modelowych obecne w~strukturach z~zewnętrzną pętlą PI: (i)~wrażliwość na~niedokładność wartości pasywnych $L$, $C$, (ii)~konieczność modelowania strat przewodzenia i~przełączania, (iii)~konieczność jawnej~estymacji prądu obciążenia (typowo realizowanej obserwatorem stanu lub~analogowym filtrem dolnoprzepustowym) oraz~(iv)~kompromisy w~doborze nastaw regulatora PI omawiane w~podrozdziale~\ref{subsec:pi}. Algorytm MF-BB osiąga ten~rezultat poprzez~włączenie do~prawa sterowania członu predykcyjnego, którego waga~$w_i$ jest~głównym parametrem strojącym -- regulator ekstrapoluje krótkoterminowy kierunek zmian napięcia wyjściowego i~na~tej~podstawie wyprzedzająco modyfikuje wartość zadaną prądu cewki, kompensując w~ten~sposób dynamikę fazy nieminimalnej (podrozdział~\ref{sec:nmp}) bez~potrzeby formalnej syntezy regulatora w~przestrzeni częstotliwości.

Dodatkową cechą podnoszącą praktyczną przydatność MF-BB w~zastosowaniach mikrosieciowych jest~niski wymagany budżet obliczeniowy -- pojedynczy krok algorytmu sprowadza się do~kilku operacji arytmetycznych (różnica, mnożenie wagą predyktora, filtracja dolnoprzepustowa, porównanie z~progiem histerezy), co~pozwala na~implementację nawet na~najprostszych mikrokontrolerach 8-bitowych pracujących z~częstotliwością MHz. Stanowi to~istotną przewagę nad~omawianym wcześniej FS-MPC, którego budżet predykcyjny rośnie eksponencjalnie z~horyzontem $N_p$ (podrozdział~\ref{sec:fsmpc}).

W~praktycznej implementacji prezentowanej w~niniejszej pracy (rozdział~\ref{ch:realizacja}), algorytm MF-BB wzbogacono dodatkowo o~mechanizm estymacji prądu obciążenia oraz~cyfrowy filtr dolnoprzepustowy w~ścieżce predyktora, których nastawy ($w_i$, $f_c$) podlegają automatycznemu strojeniu metodą metaheurystyczną PSO (rozdział~\ref{ch:ai}).

\section{Podsumowanie przeglądu}
\label{sec:przeglad-podsumowanie}

Przegląd omawiany w~niniejszym rozdziale obejmował systematyczną analizę ewolucji strategii sterowania dla~dwukierunkowych przekształtników DC-DC w~zastosowaniach mikrosieciowych, poczynając od~topologicznych uwarunkowań sprzętu (sekcja~\ref{sec:przetwornice}), poprzez~fundamentalne ograniczenia dynamiki fazy nieminimalnej (sekcja~\ref{sec:nmp}), aż~do~porównania trzech zasadniczych podejść algorytmicznych: klasycznych regulatorów PI (sekcja~\ref{sec:klasyczne}), predykcyjnego FS-MPC (sekcja~\ref{sec:fsmpc}) oraz~bezmodelowych strategii Bang-Bang (sekcja~\ref{sec:bb}).

Syntetyzując powyższe analizy, następujące cechy czynią algorytm MF-BB optymalnym wyborem w~kontekście niniejszej pracy:

\begin{enumerate}
    \item \textbf{Eliminacja potrzeby formalnego modelowania obiektu} -- w~odróżnieniu od~klasycznego PI (wymagającego identyfikacji parametrów pasywnych oraz~strat) czy~FS-MPC (wymagającego dyskretyzowanego modelu dynamiki), MF-BB operuje wyłącznie na~mierzonych wielkościach obwodowych, bezpośrednio zawierających efekty wszystkich niemodelowanych zjawisk fizycznych.
    
    \item \textbf{Zdolność do~radzenia sobie z~dynamiką NMP} -- poprzez~włączenie członu predykcyjnego z~wagą~$w_i$, algorytm kompensuje początkową odwrotną reakcję napięciową, unikając fundamentalnego ograniczenia pasma liniowych regulatorów (sekcja~\ref{sec:nmp}).
    
    \item \textbf{Niski koszt obliczeniowy} -- pojedynczy krok sterowania wymaga~kilka operacji arytmetycznych, co~umożliwia implementację na~najprostszych mikrokontrolerach pracujących z~częstotliwością MHz, w~porównaniu z~eksponencjalnie rosnącym budżetem FS-MPC przy~rozszerzeniu horyzontu predykcji.
    
    \item \textbf{Ograniczona liczba parametrów do~strojenia} -- algorytm MF-BB wymaga doboru zaledwie dwóch parametrów $(w_i, f_c)$: wagi predyktora oraz~częstotliwości granicznej filtru dolnoprzepustowego. Stałość liczby parametrów niezależnie od~zmian topologii lub~zakresu działania przekształtnika stanowi istotną zaletę w~porównaniu z~problemem doboru~$K_p, K_i$ dla~każdego punktu pracy w~klasycznej pętli PI.
    
    \item \textbf{Naturalny interfejs do~automatycznego strojenia} -- fakt, że~tylko dwa~parametry podlegają optymalizacji oraz~że~oba~są~zmiennymi rzeczywistoliczbowymi w~przestrzeni ciągłej, czyni algorytm idealnym kandydatem do~automatycznego strojenia metodą metaheurystyczną (PSO), co~omawiane będzie szczegółowo w~rozdziale~\ref{ch:ai}.
\end{enumerate}

Wybór algorytmu PSO (Particle Swarm Optimization) jako narzędzia optymalizacji parametrów MF-BB uzasadniają następujące względy:

\begin{enumerate}
    \item \textbf{Brak potrzeby gradientów} -- w~odróżnieniu od~metod gradientowych, PSO nie~wymaga analitycznego obliczenia pochodnych cząstkowych funkcji kosztu względem~$(w_i, f_c)$. Jest~to~istotne, gdyż~funkcja kosztu (wprowadzona w~rozdziale~\ref{ch:ai}) zależy od~całej trajektorii stanu wyjściowego przekształtnika -- jej~różniczkowalność nie~jest~gwarantowana ze~względu na~nieciągłości sygnału sterującego $s\in\{0,1\}$ oraz~nieliniowości modelu fizycznego.
    
    \item \textbf{Algorytm niezależny od~wymiaru przestrzeni} -- struktura PSO skaluje się liniowo z~wymiarem problemu optymalizacji, umożliwiając ewentualne rozszerzenie liczby parametrów (np.~dodanie współczynników predyktora wyższego rzędu) bez~fundamentalnych zmian algorytmu.
    
    \item \textbf{Sprawdzona skuteczność w~strojeniu kontrolerów} -- literatura tematu (m.in.~artykuły Tatariego i~Iwańskiego) uzasadnia wybór PSO dla~problemu strojenia histerezowych regulatorów energoelektronicznych, wskazując na~zbieżność do~rozwiązań porównywalnych z~klasycznymi metodami optimizacyjnymi.
    
    \item \textbf{Duża popularność i~dojrzałość implementacyjna} -- algorytm PSO jest~stosunkowo prosty do~implementacji od~podstaw (co~omawiane będzie w~rozdziale~\ref{ch:ai}), a~jednocześnie na~tyle dobrze~znany, że~nie~wymaga formalnego dowodu zbieżności czy~uzasadniania wyboru w~dyskusjach naukowych.
\end{enumerate}

Strategia pracy realizowana w~niniejszym projekcie wynika naturalnie z~powyższych analiz: (i)~implementacja bliźniaka cyfrowego MF-BB w~języku Python z~jawnym modelowaniem dynamiki przekształtnika, (ii)~strojenie parametrów $(w_i, f_c)$ za~pomocą PSO poprzez~systematyczną ewaluację pięciu konkurujących funkcji kosztu (omówionych w~rozdziale~\ref{ch:ai}), oraz~(iii)~walidacja wyników w~środowisku referencyjnym PSIM w~celu potwierdzenia transferu wiedzy z~systemu zmodelowanego. Szczegółowa implementacja oraz~wyniki eksperymentów stanowią zawartość następnych czterech rozdziałów niniejszej pracy.
